\documentclass[10pt,a4paper]{amsart}
\usepackage[margin=19mm]{geometry}
\usepackage{amsmath,amssymb,mathtools}
\usepackage[scr=rsfs,cal=euler]{mathalfa}
\usepackage{xfrac}
\usepackage[unicode,hidelinks,bookmarks=false]{hyperref}
\newcommand{\ie}{i.e.\ }
\newcommand{\cf}{cf.\ }
\newcommand{\resp}{respectively\ }
\newcommand{\Subsection}{\subsection}
\providecommand{\nobreakdash}{\nobreak\hbox{-}}
\setlength{\emergencystretch}{2em}
\multlinegap0pt
\usepackage{bm}
\usepackage{bbm}
\usepackage{dsfont}
\usepackage{xspace}
\usepackage{cancel}
\usepackage{mathtools}
\usepackage{amsthm}
\makeatletter
\@ifundefined{lemma}{\newtheorem{lemma}{Lemma}}{}
\@ifundefined{theorem}{\newtheorem{theorem}{Theorem}}{}
\makeatother
\title[Line Multigraphs of General Hypergraphs]{Line Multigraphs of General Hypergraphs}
\author{Kau\^e Cardoso}
\date{Source: arXiv:2601.09174v2 (CC BY 4.0)}
\begin{document}
\maketitle

\noindent\emph{Source and licence.} Line Multigraphs of General Hypergraphs. Version arXiv:2601.09174v2 (CC BY 4.0), distributed under the Creative Commons Attribution 4.0 International licence, as recorded in the arXiv licence metadata for this exact version and shown on the abstract page https://arxiv.org/abs/2601.09174v2. Full rights notice: licenses/arxiv-2601.09174-CC-BY-4.0.txt.\par\smallskip

\noindent\textit{Editorial scope.} Counted unit: the Theorem whose source label is \texttt{teo:collar\_eigenvalue} in the article (source0.tex lines 458--460, source label \texttt{teo:collar\_eigenvalue}), delivered together with the article's own complete original proof of it (source0.tex lines 462--479, one \texttt{proof} environment of 18 lines). No mathematics is written, repaired or re-derived: statement and proof are the article's own text line for line. Required context retained uncounted: lem:lk (lines 431--438). Every definition, display and label of those lines is retained unchanged. Omitted from this package and not redistributed: 1--430, 439--457, 461--461, 480--705. Nothing inside a retained excerpt is omitted or altered: every retained body is delivered exactly as the article's lines are written, so no omission receipt of any kind accompanies this package. Citations: the retained excerpts carry no citation command at all, so no bibliography entry of the article is transcribed and no reference is redistributed. Reference adaptation: none was needed and none was made; every reference inside the delivered unit resolves inside it, so no unresolved reference or citation is produced. Printed slips kept literal: no printed word, symbol, hypothesis or number of the counted statement or of its proof was altered. Layout and class: the article's own class is \texttt{amsart} with packages including geometry, amssymb, amsbsy, epsfig, verbatim, ulem, hyperref, tikz; the wrapper is the lane's pinned \texttt{amsart} editorial wrapper at 10pt on A4, which restates the theorem-like environments the retained text needs and adds the article's own macro definitions that the retained text needs (none), with extra wrapper packages (bm, bbm, dsfont, xspace, cancel, mathtools). The delivered numbering is the unit's own. Assets: no retained span contains a figure environment, a graphics-inclusion command, a PDF-page inclusion, a table or a TikZ picture; no image file is copied into this package, the package declares no asset dependency and no image file is redistributed with it. Licence: version arXiv:2601.09174v2 of this article is distributed under the Creative Commons Attribution 4.0 International licence, as recorded in the arXiv licence metadata for this exact version and shown on the abstract page https://arxiv.org/abs/2601.09174v2. A full rights notice is delivered at licenses/arxiv-2601.09174-CC-BY-4.0.txt. Characters: every retained excerpt is pure ASCII, so the delivered engine, XeLaTeX with its default Latin Modern text font, reports no missing character.
\begin{lemma}\label{lem:lk}
Let $\mathcal{H}$ be a hypergraph with rank $r$ and $m$ edges. The value $-r$ is an eigenvalue of the adjacency matrix of the line multigraph $\mathcal{L}(\mathcal{H})$ if and only if there exists a non-zero vector $\mathbf{x} \in \mathbb{R}^{m}$ such that:
\begin{itemize}
\item $\mathbf{B}(\mathcal{H})\mathbf{x} = \mathbf{0}$;

\item $x_i = 0$ for every hyperedge $e_i \in E(\mathcal{H})$ with cardinality $|e_i| < r$.
\end{itemize}
\end{lemma}

\begin{theorem}\label{teo:collar_eigenvalue}
Let $\mathcal{H}$ be a hypergraph with rank $r$ and $m$ hyperedges. If $\mathcal{H}$ contains an $r$-uniform collar $\mathcal{C}$ as a sub-hypergraph, then $\lambda = -r$ is an eigenvalue of the adjacency matrix $\mathbf{A}_{\mathcal{L}}$ of its line multigraph.
\end{theorem}

\begin{proof}
To prove that $-r$ is an eigenvalue, by Lemma \ref{lem:lk}, it suffices to exhibit a non-zero vector $\mathbf{x} \in \mathbb{R}^{m}$ with $\mathbf{B}\mathbf{x} = \mathbf{0}$ and $x_i = 0$ for every hyperedge $e_i$ with $|e_i| < r$.

Let $E(\mathcal{C}) \subseteq E(\mathcal{H})$ be the set of hyperedges in the collar sub-hypergraph; since $\mathcal{C}$ is $r$-uniform, every hyperedge in $E(\mathcal{C})$ has cardinality $r$, so $E(\mathcal{C}) \subseteq E_{[r]}$. Since $\mathcal{C}$ is a collar, it admits an edge $2$-coloring $f: E(\mathcal{C}) \to \{1, 2\}$. We define the vector $\mathbf{x} \in \mathbb{R}^m$ as:
$$x_i = \begin{cases} 
\;\;\,1 & \text{if } e_i \in E(\mathcal{C}) \text{ and } f(e_i) = 1, \\ 
-1 & \text{if } e_i \in E(\mathcal{C}) \text{ and } f(e_i) = 2, \\
\;\;\,0 & \text{if } e_i \notin E(\mathcal{C}).
\end{cases}$$
Since $x_i \neq 0$ only for $e_i \in E(\mathcal{C})$, and all such hyperedges have cardinality $r$, the second condition of Lemma \ref{lem:lk} holds.

Now, we examine the entries of the product $\mathbf{B}\mathbf{x}$ for each vertex $v \in V(\mathcal{H})$:
\begin{itemize}
\item \textbf{Case 1 ($v \notin V(\mathcal{C})$):} If the vertex $v$ is not part of the collar, it is not incident to any hyperedge $e_i$ for which $x_i \neq 0$. Thus, $(\mathbf{B}\mathbf{x})_v = 0$.
\item \textbf{Case 2 ($v \in V(\mathcal{C})$):} Since $\mathcal{C}$ is a $2$-regular sub-hypergraph, $v$ is incident to exactly two hyperedges in $E(\mathcal{C})$. By the definition of the edge $2$-coloring of a collar, one of these hyperedges must have color $1$ ($x_i = 1$) and the other must have color $2$ ($x_j = -1$). Consequently, $(\mathbf{B}\mathbf{x})_v = 1 + (-1) = 0.$
\end{itemize}
In both cases, we find that $\mathbf{B}\mathbf{x} = \mathbf{0}$. Since $\mathbf{x}$ is a non-zero vector satisfying both conditions of Lemma \ref{lem:lk}, it follows that $\lambda = -r$ is an eigenvalue of $\mathbf{A}_{\mathcal{L}}$.
\end{proof}

\end{document}
